\newpage
\section{Misión 9: Pasillo interplanetario}\label{sec:mission9}

\begin{minipage}{0.2\linewidth}
% \vspace{-0.5cm}
\includegraphics[width=\linewidth]{hallway}
\end{minipage}
\hspace{0.03\linewidth}
\begin{minipage}{0.80\linewidth}
Gracias en parte a tu ayuda configurando la red y los servicios de la colonia,
tu colonia florece con actividad cultural y científica. El pasillo principal del puerto espacial interplanetario de tu colonia recibe cada día a miles de visitantes de distintos planetas.
Estos visitantes exigen conectividad a Internet rápida y sencilla desde este pasillo,
y la asamblea de la colonia te ha designado para ayudar con la configuración.
\end{minipage}\ \\[0.2cm]

\subsection*{Detalles del informe}

La asamblea te ha pedido que configures un servidor DHCP para los dispositivos de los pasajeros.
El Sistema Autónomo número 7847 te presta amablemente las $1024$ direcciones más altas de su bloque, que serán las proporcionadas por este servicio.\\

\noindent
Cada $2$ horas llega al puerto espacial una nave con hasta $300$ personas (cada una con un único dispositivo ávido de Internet). Algunas se marchan de inmediato mientras otras permanecen más tiempo en el puerto espacial, pero como mucho $1000$
dispositivos necesitan estar conectados al mismo tiempo.
%
Justo antes de que el puerto espacial comience a operar por la mañana, todas las direcciones gestionadas por el servidor DHCP deben estar disponibles para su asignación, porque las $5$ primeras naves vienen completamente llenas.

\begin{enumerate}[itemsep=0.1cm]
\item Averigua quién te está prestando generosamente las $1024$ direcciones IP públicas, y cuál es su bloque completo.

\item Proporciona una tabla de encaminamiento actualizada para el router (encaminador) $\textrm{R}_0$. Este es el router conectado directamente a la \mbox{HugeAntenna\texttrademark} de la colonia, a los buses Norte y Sur, y ahora a la nueva LAN (red de área local) dedicada a los pasajeros. Puedes suponer que esta nueva LAN está conectada a una interfaz recién añadida.

\item ¿Cuál es el Lease Time más largo que puedes configurar que sea compatible con la avalancha matinal? Justifica tu respuesta.

\item Usando diagramas swimlane, proporciona un escenario concreto en el que el máximo de $1000$ dispositivos esté conectado simultáneamente en algún momento durante la hora punta. Indica cuántos mensajes de cada subtipo DHCP intervienen.

\item Si un pasajero se escondiera en la terminal del puerto espacial durante la noche, ¿qué mecanismo usarías para asegurar que su dirección IP asignada esté disponible para su asignación a la mañana siguiente?

\end{enumerate}


\subsection*{Detalles de la retroalimentación}

Considera requisitos y un proceso similares a los de las misiones anteriores, asegurándote de tratar los siguientes aspectos clave:
\begin{itemize}
\item ¿Es correcta la tabla de encaminamiento, y tan simple como sea posible?
\item ¿Has encontrado un Lease Time más largo que sea compatible con los requisitos?
\item ¿Ves algún problema en la descripción del escenario de $1000$ dispositivos concurrentes?
\end{itemize}

